\documentclass[12pt,a4paper]{article}
\usepackage[utf8]{inputenc}
\usepackage[T1]{fontenc}
\usepackage[brazil]{babel}
\usepackage{listings}
\usepackage{geometry}
\geometry{margin=2.5cm}

\title{Laboratório de Mensageria Pub/Sub com AMQP}
\author{BCC}
\date{\today}

\begin{document}
\maketitle

\section{Objetivo}
Este laboratório demonstra comunicação assíncrona com Apache Qpid JMS e um broker AMQP local.
Foram implementados produtores e consumidores para filas, publicação/assinatura em tópicos e entrega persistente.

\section{Experimentos}
\begin{enumerate}
  \item \textbf{Fila}: mensagens são enviadas por \texttt{Produtor} e recebidas por \texttt{Consumidor}.
  \item \textbf{Fan-out}: \texttt{PublicadorTopico} publica uma mensagem e cada \texttt{AssinanteTopico} ativo recebe uma cópia.
  \item \textbf{Persistência}: \texttt{ProdutorPersistente} usa \texttt{DeliveryMode.PERSISTENT}; a mensagem pode ser consumida após a reconexão.
\end{enumerate}

\section{Resultados}
Os registros em \texttt{docs/prints/} mostram mensagens enviadas, recebidas e publicadas.
Em uma execução com o broker em \texttt{localhost:61616}, os três cenários foram representados pelos respectivos programas Java.

\section{Execução}
Após iniciar o broker, execute \texttt{mvn compile} e rode as classes com a IDE ou com o classpath gerado pelo Maven.

\end{document}
